\documentclass[10.5pt,a4paper]{article}
% Estilo visual Ryoku: papel y tinta (PAPER #000000, INK #cdc4ba), sin sombras, radio minimo
\usepackage[utf8]{inputenc}
\usepackage[T1]{fontenc}
\usepackage{XCharter}          % titulares (alternativa serif a Fraunces en pdflatex)
\usepackage[sfdefault,scaled=0.95]{FiraSans}  % UI y cuerpo (Inter y Space Grotesk no estan en pdflatex)
\usepackage[scaled=0.9]{inconsolata}   % datos y etiquetas (alternativa a Space Mono)
\renewcommand{\familydefault}{\sfdefault}
\usepackage[spanish]{babel}
\usepackage[margin=2.3cm,top=2.2cm,bottom=2.4cm]{geometry}
\usepackage[table]{xcolor}
\usepackage{microtype,booktabs,tabularx,longtable,tikz,titlesec,fancyhdr,enumitem}
\usepackage[most]{tcolorbox}
\usepackage{needspace}
\newcommand{\cab}[1]{\mbox{\mono\footnotesize\color{inkmuted}#1}}
\definecolor{paper}{HTML}{000000}
\definecolor{ink}{HTML}{CDC4BA}
\definecolor{inkdim}{HTML}{B0A9A0}
\definecolor{inkmuted}{HTML}{958F87}
\definecolor{tint}{HTML}{121110}
\definecolor{tinthi}{HTML}{1E1C1A}
\definecolor{inkline}{HTML}{3A3733}   % ink al ~.26 sobre negro
\pagecolor{paper}\color{ink}
\usepackage[colorlinks=true,linkcolor=ink,urlcolor=inkdim,citecolor=ink]{hyperref}

\newcommand{\serif}{\rmfamily}
\newcommand{\mono}{\ttfamily}
\newcommand{\eyebrow}[1]{{\mono\footnotesize\color{inkmuted}\MakeUppercase{#1}}}
\newcommand{\cSi}{\textcolor{ink}{\large$\bullet$}}\newcommand{\cDudo}{\textcolor{inkmuted}{?}}\newcommand{\cNada}{\textcolor{inkmuted}{--}}
\newcommand{\nv}{\textcolor{inkmuted}{\textit{No verificado}}}

\titleformat{\section}{\serif\huge\color{ink}}{\mono\small\color{inkmuted}\twodigits{\thesection}\ //\ }{0pt}{}[\vspace{2pt}{\color{inkline}\titlerule}]
\titleformat{\subsection}[hang]{\serif\fontsize{28}{30}\selectfont\color{ink}}{\mono\small\color{inkmuted}\thesubsection\ //\ }{0pt}{}
\titleformat{name=\subsection,numberless}{\mono\small\color{inkdim}}{}{0pt}{}\titlespacing*{\subsection}{0pt}{16pt}{6pt}
\titlespacing*{\section}{0pt}{22pt}{10pt}
\titlespacing*{\subsection}{0pt}{16pt}{6pt}
\newcommand{\twodigits}[1]{\ifnum#1<10 0#1\else#1\fi}

\pagestyle{fancy}\fancyhf{}
\renewcommand{\headrulewidth}{0pt}
\fancyfoot[L]{\mono\footnotesize\color{inkmuted}RYOKU \textperiodcentered\ \docid}
\fancyfoot[R]{\mono\footnotesize\color{inkmuted}\thepage}
\fancypagestyle{plain}{\fancyhf{}\fancyfoot[R]{\mono\footnotesize\color{inkmuted}\thepage}}
\arrayrulecolor{inkline}
\setlength{\aboverulesep}{3pt}\setlength{\belowrulesep}{3pt}
\setlength{\parindent}{0pt}\setlength{\parskip}{5pt}
\setlist[itemize]{leftmargin=1.2em,itemsep=3pt,label={\color{inkmuted}--}}

\newcommand{\etiq}[1]{\tikz[baseline=(t.base)]\node[draw=inkmuted,line width=0.4pt,rounded corners=2pt,inner xsep=6pt,inner ysep=3pt,text=ink,font=\mono\footnotesize] (t) {\MakeUppercase{#1}};}
\newcommand{\kpi}[2]{\begin{minipage}[t]{0.235\textwidth}{\mono\scriptsize\color{inkmuted}\MakeUppercase{#1}}\\[2pt]{\serif\large\color{ink}#2}\end{minipage}}
\newenvironment{ficha}[6]{\needspace{22\baselineskip}\subsection{#1}\par\nopagebreak\vspace{-4pt}\etiq{#2}\par\vspace{8pt}%
 \begin{tcolorbox}[enhanced,colback=tint,colframe=inkline,boxrule=0.4pt,arc=2pt,left=8pt,right=8pt,top=7pt,bottom=7pt,boxsep=0pt]
 \kpi{Empleados}{#3}\hfill\kpi{Fundaci\'on}{#4}\hfill\kpi{Modelo}{#5}\hfill\kpi{Cloud}{#6}
 \end{tcolorbox}\vspace{6pt}{\color{inkline}\hrule height 0.4pt}\vspace{3pt}}{}
\newcommand{\campo}[2]{\noindent\hangindent=3.4cm\hangafter=1\makebox[3.4cm][l]{\mono\footnotesize\color{inkmuted}\MakeUppercase{#1}}#2\par\vspace{2pt}{\color{inkline}\hrule height 0.4pt}\vspace{3pt}}

\newcommand{\cPar}{\textcolor{inkmuted}{$\circ$}}
\newcommand{\crit}{\textcolor{ink}{\large$\bullet$}}
\newcommand{\perfilkpi}[2]{\begin{minipage}[t]{0.235\textwidth}{\mono\scriptsize\color{inkmuted}\MakeUppercase{#1}}\\[2pt]{\serif\large\color{ink}#2}\end{minipage}}

\newcommand{\docid}{RA1b}
\tikzset{
  org/.style={draw=inkmuted,line width=0.4pt,rounded corners=2pt,fill=tinthi,text=ink,font=\sffamily\footnotesize,align=center,minimum height=0.9cm,inner sep=3pt},
  crit/.style={fill=ink,draw=ink,text=paper,font=\sffamily\footnotesize\bfseries},
  cima/.style={org,minimum width=3.6cm},
  leaf/.style={org,fill=tint,draw=inkline,font=\sffamily\scriptsize,minimum height=0.6cm},
  lk/.style={draw=inkmuted,line width=0.4pt},
}
\newcommand{\leyenda}{\vspace{4pt}{\mono\footnotesize\color{inkmuted}%
 \tikz[baseline=-0.5ex]\node[org,crit,minimum height=0.35cm,minimum width=0.7cm]{};\ departamento cr\'itico\quad
 \tikz[baseline=-0.5ex]\node[org,minimum height=0.35cm,minimum width=0.7cm]{};\ departamento\quad
 \tikz[baseline=-0.5ex]\node[leaf,minimum height=0.35cm,minimum width=0.7cm]{};\ equipo o funci\'on}\par}

\begin{document}
\thispagestyle{empty}
\begin{tikzpicture}[remember picture,overlay]
  \fill[ink,opacity=0.16] ([xshift=-3.2cm,yshift=-7.2cm]current page.north east) circle (4.6cm);
  \draw[inkline,line width=0.4pt] ([xshift=2.3cm,yshift=-1.6cm]current page.north west) -- ([xshift=-2.3cm,yshift=-1.6cm]current page.north east);
  \node[anchor=north west,inner sep=0] at ([xshift=2.3cm,yshift=-1.0cm]current page.north west) {\eyebrow{00 // Proyecto intermodular DAM}};
  \node[anchor=north east,inner sep=0] at ([xshift=-2.3cm,yshift=-1.0cm]current page.north east) {\eyebrow{2026/2027}};
  \node[anchor=south west,inner sep=0,text width=14cm] at ([xshift=2.3cm,yshift=7.5cm]current page.south west)
    {{\serif\fontsize{54}{56}\selectfont\color{ink} Mapa\\ organizativo\\ sectorial\par}};
  \node[anchor=south west,inner sep=0,text width=14cm] at ([xshift=2.3cm,yshift=5.2cm]current page.south west)
    {{\mono\color{inkdim}RA1b \textperiodcentered\ 3 organigramas tipo \textperiodcentered\ 14 oct 2026\par}};
  \draw[inkline,line width=0.4pt] ([xshift=2.3cm,yshift=3.2cm]current page.south west) -- ([xshift=-2.3cm,yshift=3.2cm]current page.south east);
  \node[anchor=north west,inner sep=0] at ([xshift=2.3cm,yshift=3.0cm]current page.south west)
    {\eyebrow{PAPER \#000000 \textperiodcentered\ INK \#cdc4ba \textperiodcentered\ GRAIN 5.5\% \textperiodcentered\ RADIUS 2 \textperiodcentered\ NO SHADOW}};
\end{tikzpicture}
\newpage
\setcounter{tocdepth}{2}
\renewcommand{\contentsname}{\mono\small\color{inkmuted}\MakeUppercase{\'Indice}}
\tableofcontents
\newpage

\section{Objetivo y metodolog\'ia}
Se dise\~nan tres organigramas tipo para una empresa de software (startup, pyme y gran empresa) y se identifican los departamentos cr\'iticos de cada una. Un departamento es \textbf{cr\'itico} cuando, si falla o no existe, la empresa deja de entregar su producto o de ingresar dinero.

Los organigramas son \textbf{modelos gen\'ericos}, no el de una empresa concreta. Parten de lo observado en RA1a (Factorial y Holded como startups, Autentia y Plain Concepts como pymes, Indra y Amadeus como grandes) y de criterio propio. Los tama\~nos son orientativos y los nombres de los cargos var\'ian entre empresas. La clasificaci\'on de lo que es cr\'itico es una valoraci\'on propia, no un dato de una fuente.

\section{Resumen comparativo}
\rowcolors{2}{tint}{paper}
\begin{center}\small
\begin{tabularx}{\textwidth}{@{}>{\bfseries}l >{\raggedright\arraybackslash}p{2.6cm} >{\raggedright\arraybackslash}p{2.9cm} X@{}}
\rowcolor{paper}\cab{TIPO} & \cab{TAMA\~NO ORIENTATIVO} & \cab{ESTRUCTURA} & \cab{DEPARTAMENTOS CR\'ITICOS} \\
\arrayrulecolor{inkmuted}\midrule
Startup & 2--50 personas & plana, 2 niveles & Ingenier\'ia, Producto, Ventas y crecimiento \\
Pyme & 50--250 personas & funcional, 3 niveles & Desarrollo, Proyectos y entrega, Comercial \\
Gran empresa & 1.000+ personas & por divisiones, 4--5 niveles & Ingenier\'ia, Ciberseguridad, Datos, Infraestructura y cloud, Comercial \\
\arrayrulecolor{inkmuted}\bottomrule
\end{tabularx}
\end{center}
\rowcolors{1}{paper}{paper}

\clearpage
\section{Organigramas tipo}

\subsection{Startup}
\etiq{2--50 personas \textperiodcentered\ producto propio}\par\vspace{6pt}
\begin{center}
\begin{tikzpicture}
  \node[org,cima] (ceo) at (0,0) {Fundador/a \textperiodcentered\ CEO};
  \node[org,crit] (cto) at (-5.6,-1.9) {CTO\\Ingenier\'ia};
  \node[org,crit] (pro) at (-1.9,-1.9) {Producto\\y dise\~no};
  \node[org,crit] (gro) at (1.9,-1.9) {Ventas y\\crecimiento};
  \node[org] (ope) at (5.6,-1.9) {Operaciones\\y administraci\'on};
  \draw[lk] (ceo.south) -- ++(0,-0.45) coordinate (b);
  \draw[lk] (-5.6,0 |- b) -- (5.6,0 |- b);
  \foreach \x in {-5.6,-1.9,1.9,5.6} \draw[lk] (\x,0 |- b) -- (\x,-1.5);
  \node[leaf,minimum width=3.1cm] (l1) at (-5.6,-3.2) {Backend};
  \node[leaf,minimum width=3.1cm] (l2) at (-5.6,-3.95) {Frontend / m\'ovil};
  \node[leaf,minimum width=3.1cm] (l3) at (-5.6,-4.7) {DevOps (compartido)};
  \node[leaf,minimum width=3.1cm] (l4) at (-1.9,-3.2) {Producto (PM)};
  \node[leaf,minimum width=3.1cm] (l5) at (-1.9,-3.95) {UX / UI};
  \node[leaf,minimum width=3.1cm] (l6) at (1.9,-3.2) {Ventas};
  \node[leaf,minimum width=3.1cm] (l7) at (1.9,-3.95) {Marketing};
  \node[leaf,minimum width=3.1cm] (l8) at (1.9,-4.7) {Atenci\'on al cliente};
  \node[leaf,minimum width=3.1cm] (l9) at (5.6,-3.2) {Finanzas y legal};
  \node[leaf,minimum width=3.1cm] (l10) at (5.6,-3.95) {Personas (RR.\,HH.)};
  \foreach \p/\q in {cto/l1,pro/l4,gro/l6,ope/l9} \draw[lk] (\p.south) -- (\q.north);
  \foreach \p/\q in {l1/l2,l2/l3,l4/l5,l6/l7,l7/l8,l9/l10} \draw[lk] (\p.south) -- (\q.north);
\end{tikzpicture}
\end{center}
\leyenda
\subsection*{Departamentos cr\'iticos}
\campo{Ingenier\'ia}{es el producto. Con un equipo de 3 a 10 personas, una baja o un mal fichaje frena el lanzamiento. DevOps no es un equipo aparte: lo asume el propio equipo de desarrollo}
\campo{Producto y dise\~no}{decide qu\'e se construye. Construir lo equivocado agota la caja antes de encontrar el mercado}
\campo{Ventas y crecimiento}{sin clientes de pago no hay siguiente ronda de financiaci\'on. En SaaS (Factorial, Holded en RA1a) es el motor de ingresos}
\campo{Operaciones}{no es cr\'itico: finanzas, legal y personas las lleva el fundador o un gestor externo hasta que haya volumen}

\clearpage
\subsection{Pyme}
\etiq{50--250 personas \textperiodcentered\ consultor\'ia o producto}\par\vspace{6pt}
\begin{center}
\begin{tikzpicture}
  \node[org,cima] (dg) at (0,0) {Direcci\'on general};
  \node[org] (dt) at (-5.6,-1.9) {Direcci\'on\\t\'ecnica};
  \node[org,crit] (pm) at (-1.9,-1.9) {Proyectos\\y entrega};
  \node[org,crit] (dc) at (1.9,-1.9) {Direcci\'on\\comercial};
  \node[org] (da) at (5.6,-1.9) {Administraci\'on\\y personas};
  \draw[lk] (dg.south) -- ++(0,-0.45) coordinate (b);
  \draw[lk] (-5.6,0 |- b) -- (5.6,0 |- b);
  \foreach \x in {-5.6,-1.9,1.9,5.6} \draw[lk] (\x,0 |- b) -- (\x,-1.5);
  \node[leaf,crit,minimum width=3.1cm] (a1) at (-5.6,-3.2) {Desarrollo};
  \node[leaf,minimum width=3.1cm] (a2) at (-5.6,-3.95) {QA y pruebas};
  \node[leaf,minimum width=3.1cm] (a3) at (-5.6,-4.7) {Sistemas e infraestructura};
  \node[leaf,minimum width=3.1cm] (a4) at (-5.6,-5.45) {Arquitectura};
  \node[leaf,minimum width=3.1cm] (b1) at (-1.9,-3.2) {Jefes de proyecto};
  \node[leaf,minimum width=3.1cm] (b2) at (-1.9,-3.95) {An\'alisis funcional};
  \node[leaf,minimum width=3.1cm] (b3) at (-1.9,-4.7) {Equipos de proyecto};
  \node[leaf,minimum width=3.1cm] (c1) at (1.9,-3.2) {Preventa};
  \node[leaf,minimum width=3.1cm] (c2) at (1.9,-3.95) {Cuentas y ventas};
  \node[leaf,minimum width=3.1cm] (c3) at (1.9,-4.7) {Marketing};
  \node[leaf,minimum width=3.1cm] (d1) at (5.6,-3.2) {Finanzas};
  \node[leaf,minimum width=3.1cm] (d2) at (5.6,-3.95) {RR.\,HH. y selecci\'on};
  \node[leaf,minimum width=3.1cm] (d3) at (5.6,-4.7) {Legal y calidad};
  \foreach \p/\q in {dt/a1,pm/b1,dc/c1,da/d1} \draw[lk] (\p.south) -- (\q.north);
  \foreach \p/\q in {a1/a2,a2/a3,a3/a4,b1/b2,b2/b3,c1/c2,c2/c3,d1/d2,d2/d3} \draw[lk] (\p.south) -- (\q.north);
\end{tikzpicture}
\end{center}
\leyenda
\subsection*{Departamentos cr\'iticos}
\campo{Desarrollo}{es lo que se factura. En una consultora (Autentia, Plain Concepts en RA1a) las horas del equipo son el ingreso}
\campo{Proyectos y entrega}{coordina plazos, alcance y cliente. Si falla, los proyectos se retrasan o dan p\'erdidas aunque el c\'odigo sea bueno}
\campo{Comercial}{la pyme depende de pocos clientes grandes; perder uno afecta a la ocupaci\'on del equipo}
\campo{Resto}{necesarios pero no cr\'iticos: pueden tener pocas personas o externalizarse parcialmente}

\clearpage
\subsection{Gran empresa}
\etiq{1.000+ personas \textperiodcentered\ divisiones}\par\vspace{6pt}
\begin{center}
\begin{tikzpicture}
  \node[org,cima] (ceo) at (0,0) {Consejo \textperiodcentered\ CEO};
  \def\w{2.4}
  \node[org,crit,minimum width=\w cm] (t) at (-6.4,-1.9) {CTO\\Tecnolog\'ia};
  \node[org,crit,minimum width=\w cm] (s) at (-3.85,-1.9) {CISO\\Seguridad};
  \node[org,crit,minimum width=\w cm] (dd) at (-1.3,-1.9) {Datos\\e IA};
  \node[org,crit,minimum width=\w cm] (c) at (1.3,-1.9) {CCO\\Comercial};
  \node[org,minimum width=\w cm] (o) at (3.85,-1.9) {COO\\Operaciones};
  \node[org,minimum width=\w cm] (f) at (6.4,-1.9) {CFO\\Finanzas};
  \draw[lk] (ceo.south) -- ++(0,-0.45) coordinate (b);
  \draw[lk] (-6.4,0 |- b) -- (6.4,0 |- b);
  \foreach \x in {-6.4,-3.85,-1.3,1.3,3.85,6.4} \draw[lk] (\x,0 |- b) -- (\x,-1.5);
  \node[leaf,minimum width=\w cm] (t1) at (-6.4,-3.2) {Ingenier\'ia};
  \node[leaf,minimum width=\w cm] (t2) at (-6.4,-3.95) {Arquitectura};
  \node[leaf,minimum width=\w cm] (t3) at (-6.4,-4.7) {Infraestructura y cloud};
  \node[leaf,minimum width=\w cm] (t4) at (-6.4,-5.45) {DevOps / SRE};
  \node[leaf,minimum width=\w cm] (t5) at (-6.4,-6.2) {QA};
  \node[leaf,minimum width=\w cm] (s1) at (-3.85,-3.2) {Seguridad (SOC)};
  \node[leaf,minimum width=\w cm] (s2) at (-3.85,-3.95) {Cumplimiento};
  \node[leaf,minimum width=\w cm] (d1) at (-1.3,-3.2) {Ingenier\'ia de datos};
  \node[leaf,minimum width=\w cm] (d2) at (-1.3,-3.95) {An\'alisis y BI};
  \node[leaf,minimum width=\w cm] (d3) at (-1.3,-4.7) {IA / ML};
  \node[leaf,minimum width=\w cm] (c1) at (1.3,-3.2) {Ventas};
  \node[leaf,minimum width=\w cm] (c2) at (1.3,-3.95) {Marketing};
  \node[leaf,minimum width=\w cm] (c3) at (1.3,-4.7) {\'Exito de cliente};
  \node[leaf,minimum width=\w cm] (o1) at (3.85,-3.2) {Entrega de proyectos};
  \node[leaf,minimum width=\w cm] (o2) at (3.85,-3.95) {Soporte};
  \node[leaf,minimum width=\w cm] (o3) at (3.85,-4.7) {Compras};
  \node[leaf,minimum width=\w cm] (f1) at (6.4,-3.2) {Contabilidad};
  \node[leaf,minimum width=\w cm] (f2) at (6.4,-3.95) {Legal};
  \node[leaf,minimum width=\w cm] (f3) at (6.4,-4.7) {Personas (RR.\,HH.)};
  \foreach \p/\q in {t/t1,s/s1,dd/d1,c/c1,o/o1,f/f1} \draw[lk] (\p.south) -- (\q.north);
  \foreach \p/\q in {t1/t2,t2/t3,t3/t4,t4/t5,s1/s2,d1/d2,d2/d3,c1/c2,c2/c3,o1/o2,o2/o3,f1/f2,f2/f3} \draw[lk] (\p.south) -- (\q.north);
\end{tikzpicture}
\end{center}
\leyenda
\subsection*{Departamentos cr\'iticos}
\campo{Tecnolog\'ia}{ingenier\'ia, arquitectura e infraestructura sostienen el servicio. Una ca\'ida en cloud o en producci\'on se mide en dinero por minuto}
\campo{Seguridad}{a esta escala hay datos de muchos clientes y obligaciones legales (RGPD) y de certificaci\'on (ISO 27001, vistas en RA1a). Una brecha da\~na la reputaci\'on y los contratos}
\campo{Datos e IA}{base de las decisiones y de los productos con IA; en las multinacionales de RA1a es una l\'inea de negocio propia}
\campo{Comercial}{sostiene el crecimiento y los grandes contratos}
\campo{Operaciones y finanzas}{imprescindibles, pero repetibles y sustituibles: se estandarizan y se apoyan en proveedores}

\clearpage
\section{Matriz de departamentos cr\'iticos}
\rowcolors{2}{tint}{paper}
\begin{center}\small
\begin{tabularx}{\textwidth}{@{}>{\bfseries}l *{3}{>{\centering\arraybackslash}X}@{}}
\rowcolor{paper}\cab{DEPARTAMENTO} & \cab{STARTUP} & \cab{PYME} & \cab{GRAN EMPRESA} \\
\arrayrulecolor{inkmuted}\midrule
Ingenier\'ia / Desarrollo & \crit & \crit & \crit \\
Producto y dise\~no & \crit & \cPar & \cPar \\
QA y pruebas & \cNada & \cPar & \cPar \\
Infraestructura, DevOps y cloud & \cNada & \cPar & \crit \\
Ciberseguridad & \cNada & \cNada & \crit \\
Datos e IA & \cNada & \cNada & \crit \\
Proyectos y entrega & \cNada & \crit & \cPar \\
Comercial y ventas & \crit & \crit & \crit \\
Atenci\'on al cliente y soporte & \cPar & \cPar & \cPar \\
Finanzas y administraci\'on & \cNada & \cPar & \cPar \\
RR.\,HH. y personas & \cNada & \cPar & \cPar \\
Legal y cumplimiento & \cNada & \cNada & \cPar \\
\arrayrulecolor{inkmuted}\bottomrule
\end{tabularx}
\end{center}
\rowcolors{1}{paper}{paper}
{\mono\footnotesize\color{inkmuted}$\bullet$ cr\'itico \quad $\circ$ existe, no cr\'itico \quad -- no existe como departamento propio (lo asume otro equipo o un externo).}

\section{Conclusiones}
\begin{itemize}
\item \textbf{Ingenier\'ia y comercial son cr\'iticos en los tres tama\~nos}: sin producto o sin clientes no hay empresa.
\item Al crecer, la estructura pasa de \textbf{personas polivalentes} (startup) a \textbf{funciones} (pyme) y a \textbf{divisiones} (gran empresa).
\item Cuanto m\'as grande, m\'as departamentos aparecen como cr\'iticos por riesgo y no solo por producci\'on: \textbf{seguridad} y \textbf{datos} en la gran empresa.
\item En la pyme la pieza propia es \textbf{proyectos y entrega}, porque su negocio es vender equipos y horas.
\item Para un perfil DAM, los puestos de desarrollo, QA y DevOps (ver RA1c) est\'an en el departamento cr\'itico de los tres tipos.
\end{itemize}

\section{Limitaciones}
Los organigramas son modelos simplificados de una empresa de software; en la realidad los nombres y las dependencias cambian (por ejemplo, el CISO puede depender del CTO o del CEO). Los rangos de tama\~no son aproximados y la lista de cr\'iticos es una valoraci\'on propia, no extra\'ida de una fuente. No se ha usado ning\'un organigrama real publicado de las empresas de RA1a.

\end{document}
